%!TEX TS-program = xelatex
% vim: set fenc=utf-8
% -*- coding: UTF-8; -*-

% =============================================================
%  作者简历（规范 1.6 / 附件 6）
%  本文件由 cugthesis.cls 的 \makeauthorcv 自动载入，内容请替换成你自己的。
%
%  版式（依附件 6 实测，由模板统一控制，这里只管内容）：
%   · 页面标题「作者简历」：三号 16 磅黑体居中
%   · 小节标题「一、基本情况」：小四 12 磅黑体
%   · 正文：小四 12 磅、固定行距 20 磅、首行缩进 2 个汉字符
%
%  \cugauthorcvsection{标题}          新增一个小节（自动编号 一、二、三…）
%  \cugauthorcvinfoitem{起}{止}{说明}  一行学历/经历
% =============================================================

\cugauthorcvsection{基本情况}
姓名：张三 \quad 性别：男 \quad 民族：汉 \quad 出生年月：1994-05-02 \quad 籍贯：山西省朔州市

\cugauthorcvinfoitem{2013.09}{2017.07}{中国地质大学（武汉）工学学士}
\cugauthorcvinfoitem{2021.09}{至今}{中国地质大学（武汉）攻读博士学位}

\cugauthorcvsection{学术论文}
    \begin{enumerate}
        \item 张三, 某某某, 某某. 东昆仑造山带晚古生代花岗岩锆石 U-Pb 年代学及其地质意义[J]. 地球科学, 2025, 50(6): 2101-2118.
        \item Cheng K D, XXX X X. Geochemistry of the Late Paleozoic granitoids in the East Kunlun Orogenic Belt[J]. Journal of Earth Science, 2025, 36(3): 512-528.
        \item \dots
    \end{enumerate}

\cugauthorcvsection{获奖、专利情况}
    \begin{enumerate}
        \item 中国地质大学（武汉）研究生学业奖学金一等奖, 2023.
        \item \dots
    \end{enumerate}

\cugauthorcvsection{研究项目}
    \begin{enumerate}
        \item 国家自然科学基金面上项目：东昆仑造山带古生代岩浆作用与构造演化, 2022--2025, 项目编号 42XXXXXX, 主要参加人;
        \item \dots
    \end{enumerate}

\endinput
